\documentclass[../../../ps_a.tex]{subfiles}

\begin{document}
    \begin{enumerate}
        \item Yについての電荷保存則より，
            
            \myspaceb%
            $\displaystyle%
                -x_1 + y_1 = 0 \;.
            $\\            
            キルヒホッフ則より，
            
            \myspaceb%
            $\displaystyle%
                \myfracc{x_1}{C_0} + \myfracc{y_1}{C_0/2} = 2V_0%
                \qquad \therefore \; x_1 = y_1
                = \uwave{\myfraca{2}{3}C_0 V_0} \;.
            $\\
            すると，
            
            \myspaceb%
            $\displaystyle%
                E_1 = \myfracc{x_1}{C_0 d} = \uwave{\myfracc{2V_0}{3d}} \;, 
            $

            \myspaceb%
            $\displaystyle%
                \phi_1 = \uwave{2V_0} \;,%
                \quad \theta_1 = \myfracc{y_1}{C_0/2} = \uwave{\myfraca{4}{3}V_0} \;.
            $            
        \item 電荷保存則より，Xの電荷は変わらない．つまり，
            
            \myspaceb%
            $\displaystyle%
                x_2 = x_1 = \uwave{\myfraca{2}{3}C_0 V_0} \;.
            $\\
            キルヒホッフ則より，
            
            \myspaceb%
            $\displaystyle%
                \myfracc{y_2}{C_0/2} = V_0%
                \qquad \therefore \; y_2 = \uwave{\myfraca{1}{2}C_0 V_0} \;.
            $\\
            すると，
            
            \myspaceb%
            $\displaystyle%
                E_2 = \myfracc{x_2}{C_0 d} = \uwave{\myfracc{2V_0}{3d}} \;,
            $

            \myspaceb%
            $\displaystyle%
                \theta_2 = \uwave{V_0} \;,%
                \quad \phi_2 = \theta_2 + \myfracc{x_2}{C_0}%
                = \uwave{\myfraca{5}{3}V_0} \;.
            $
        \item Yについての電荷保存則より，
            
            \myspaceb%
            $\displaystyle%
                -x_3 + y_3 = -x_2 + y_2%
                \qquad \therefore \; -x_3 + y_3 = -\myfraca{1}{6}C_0 V_0 \;.
            $\\
            キルヒホッフ則より，
            
            \myspaceb%
            $\displaystyle%
                \myfracc{x_3}{C_0} + \myfracc{y_3}{C_0/2} = 2V_0
            $

            \myspaceb%
            $\displaystyle%
                \therefore \; x_3 = \uwave{\myfraca{7}{9}C_0 V_0} \;,%
                \quad y_3 = \uwave{\myfraca{11}{18}C_0 V_0} \;.
            $\\
            すると，
            
            \myspaceb%
            $\displaystyle%
                E_3 = \myfracc{x_3}{C_0 d} = \uwave{\myfracc{7V_0}{9d}} \;,
            $

            \myspaceb%
            $\displaystyle%
                \phi_3 = \uwave{2V_0} \;,%
                \quad \theta_3 = \myfracc{y_3}{C_0/2} = \uwave{\myfraca{11}{9}V_0} \;.
            $
        \item X下面の電荷を$x_4$，Y下面の電荷を$y_4$とすると，Yについての電荷保存則より，
            
            \myspaceb%
            $\displaystyle%
                -x_4 + y_4 = -x_3 + y_3%
                \qquad \therefore \; -x_4 + y_4 = -\myfraca{1}{6}C_0 V_0 \;.
            $\\
            キルヒホッフ則より，
            
            \myspaceb%
            $\displaystyle%
                \myfracc{x_4}{C_0/2} + \myfracc{y_4}{C_0} = 2V_0
            $

            \myspaceb%
            $\displaystyle%
                \therefore \; x_4 = \uwave{\myfraca{13}{18}C_0 V_0} \;,%
                \quad y_4 = \uwave{\myfraca{5}{9}C_0 V_0} \;.
            $\\
            $\Delta Q$は，Xの電荷の変化量に等しいから，
            
            \myspaceb%
            $\displaystyle%
                \Delta Q = x_4 - x_3 = \uwave{-\myfraca{1}{18}C_0 V_0} \;.
            $            
    \end{enumerate}
\end{document}
